\chapter{Baseline compatta e metodologia sperimentale}
\label{chap:baseline_metodologia}

% Capitolo rivisto il 2026-11-11 (figure da figure_scripts/cap2; statistiche appaiate da
% src/etsr/evaluation/comparison_metrics.py, bootstrap sulla differenza di accuratezza).
% Fonti verificate il 2026-11-04 sul repository temporal-event-spiking-research,
% branch Predictive-Temporal-Coding, commit 5a03fec: src/etsr/models/{mini_qkformer,layers,spiking,sequence}.py,
% src/etsr/encoders/count.py, configs/{models,recipes,datasets,splits}, config_resolved.yaml e
% summary.json dei run B E0 (dvslip 20260825_211710, dvsgesture 20261009_174645,
% cifar10dvs 20261009_184314, thu_eact_50_chl 20261011_093547), hardware_profile_v4.json di B su DVS-Lip,
% audit di sviluppo in artifacts/datasets/*. Nessun dato del test ufficiale.

Ogni modifica studiata nei capitoli successivi viene misurata rispetto a un riferimento, e quella misura vale quanto la precisione con cui il riferimento è stato fissato. Questo capitolo lo descrive. Presenta i quattro dataset, il modo in cui il flusso di eventi diventa l'ingresso della rete, l'architettura della baseline compatta, il protocollo con cui i modelli vengono addestrati e confrontati e le grandezze con cui se ne misurano prestazioni e costo.

La parte sperimentale procede per passi successivi, ciascuno dei quali cambia un solo fattore. Il capitolo~\ref{chap:riallocazione_budget} ridistribuisce la capacità all'interno della baseline, a parità di ingresso, e ne ricava il modello di riferimento. Il capitolo~\ref{chap:clock_input} tiene ferma questa architettura e cambia il clock dell'input. Il capitolo~\ref{chap:memoria_relazionale}, con il clock così scelto, arricchisce la memoria temporale. Ciò che resta costante lungo questo percorso è descritto qui una volta per tutte.


% ==============================================================================
%  Dataset
% ==============================================================================

\section{Dataset}
\label{sec:dataset}

La sperimentazione si svolge su quattro benchmark neuromorfici di classificazione, scelti perché affidano all'informazione temporale ruoli molto diversi. In DVS-Lip le classi sono parole pronunciate, distinte da movimenti labiali brevi e spesso simili tra loro, e l'ordine in cui si susseguono le configurazioni della bocca è parte essenziale del segnale. DVS128 Gesture raccoglie gesti ampi e in larga parte ripetitivi, alcuni dei quali si distinguono soltanto per il verso del movimento. In CIFAR10-DVS la dinamica nasce dal modo in cui immagini statiche vengono mostrate al sensore, più che dal loro contenuto. THU-EACT-50-CHL, infine, contiene azioni umane complete, riprese a una risoluzione più alta e con durate molto variabili. Un meccanismo temporale utile in uno di questi scenari non lo è necessariamente negli altri, ed è proprio questa eterogeneità a rendere il confronto informativo. La figura~\ref{fig:esempi_dataset} ne mostra un campione per ciascuno.

\paragraph{DVS-Lip}
\label{subsec:dvslip}
È un dataset di lettura labiale a eventi, introdotto insieme alla rete MSTP~\cite{tan2022dvslip}. Raccoglie 100 parole inglesi pronunciate da 40 parlanti e rilasciate come regioni di $128\times128$ pixel attorno alla bocca, con una partizione ufficiale che assegna 30 parlanti al training e 10 al test~\cite{tan2022dvslip}. Le registrazioni durano circa un secondo e contengono pochi eventi, perché il segnale utile si concentra in un'area ristretta e in variazioni di forma molto fini. Il vocabolario include di proposito 25 coppie di parole visivamente simili. Le 50 parole che le compongono e le restanti 50 formano due gruppi, sui quali si calcolano due accuratezze parziali, Acc1 e Acc2.

\paragraph{DVS128 Gesture}
\label{subsec:dvsgesture}
È stato acquisito con un sensore DVS128 e comprende 11 classi, dieci gesti di mani e braccia e una classe residua che raccoglie movimenti diversi~\cite{amir2017dvsgesture}. Le registrazioni coinvolgono 29 soggetti in diverse condizioni di illuminazione, e la partizione ufficiale riserva al test gli ultimi sei. Ogni registrazione contiene più gesti consecutivi, delimitati da annotazioni temporali, e in questo lavoro un campione coincide con un segmento annotato. I segmenti durano in mediana più di sei secondi, ma alcuni superano i diciotto. Diverse classi si distinguono soltanto per il braccio impiegato o per il verso di rotazione, e per questo trasformazioni come il ribaltamento orizzontale, che le scambierebbero tra loro, sono escluse.

\paragraph{CIFAR10-DVS}
\label{subsec:cifar10dvs}
Converte in eventi $10\,000$ immagini di CIFAR-10, mille per classe~\cite{li2017cifar10dvs}. Le immagini, mostrate su un monitor, vengono fatte muovere con un movimento regolare e ripetuto davanti a un sensore DVS128, così che i loro contorni generino eventi~\cite{li2017cifar10dvs}. Ne risultano registrazioni di durata quasi uniforme, poco più di un secondo, e le più dense del gruppo. Poiché il contenuto è statico, la struttura temporale riflette soprattutto il movimento imposto durante l'acquisizione, e il dataset fa da termine di confronto in cui è lecito attendersi che la memoria temporale pesi meno. Non esiste una partizione ufficiale, e in questo lavoro i primi cento campioni di ciascuna classe formano un insieme riservato, escluso dallo sviluppo.

\paragraph{THU-EACT-50-CHL}
\label{subsec:thueact}
THU-EACT-50 è un benchmark per il riconoscimento di 50 azioni umane con event camera~\cite{gao2023action}. Questo lavoro ne adotta la versione THU-EACT-50-CHL, acquisita con un sensore DAVIS346 a $346\times260$ pixel e composta da $2\,330$ registrazioni di 18 soggetti, che gli autori ripartiscono tra training e test senza soggetti in comune. Le durate hanno una coda lunga, da una mediana inferiore ai tre secondi fino a quasi venti, e il numero di eventi varia da circa diecimila a oltre quattro milioni per registrazione. La risoluzione maggiore moltiplica inoltre il numero di posizioni spaziali, e quindi di token, che la rete elabora a ogni passo.

\begin{figure}[t]
    \centering
    \includegraphics[width=\textwidth]{images/cap2/cap2_dataset.pdf}
    \caption[Un campione per ciascun dataset]{Un campione di sviluppo per ciascun dataset, scelto con una durata vicina alla mediana. A sinistra, tre bin della rappresentazione a finestra fissa della sezione~\ref{sec:input_baseline}, in rosso dove prevalgono gli eventi ON e in blu dove prevalgono gli OFF. A destra, la frequenza degli eventi sull'intera finestra, con le bande scure in corrispondenza dei bin mostrati e in grigio chiaro la parte successiva all'ultimo evento. I campioni sono la parola \emph{warning} per DVS-Lip, un saluto con la mano destra per DVS128 Gesture e un aereo per CIFAR10-DVS.}
    \label{fig:esempi_dataset}
\end{figure}

\paragraph{Partizioni di sviluppo}
\label{subsec:partizioni_sviluppo}

I test ufficiali non partecipano allo sviluppo. Ogni decisione, dalla configurazione dell'input alla scelta del checkpoint, si basa su una validazione ricavata dal solo training ufficiale, chiamata nel seguito validazione di sviluppo (\emph{development validation}). Le partizioni sono fissate prima di qualsiasi esperimento e sono le stesse per tutti i modelli, e ogni confronto avviene quindi sugli stessi campioni. La tabella~\ref{tab:dataset_sintesi} ne riporta le dimensioni insieme alle durate delle registrazioni.

\begin{table}[t]
\centering
\small
\renewcommand{\arraystretch}{1.2}
\caption[Dataset e partizioni di sviluppo]{Dataset e partizioni di sviluppo. La partizione riservata è il test ufficiale, tranne che per CIFAR10-DVS, per il quale è l'insieme escluso dallo sviluppo. Le durate si riferiscono alle partizioni di sviluppo e riportano il valore mediano e, tra parentesi, il massimo.}
\label{tab:dataset_sintesi}
\setlength{\tabcolsep}{4pt}
\begin{tabular}{@{}l c c r r r c@{}}
\hline
\textbf{Dataset} & \textbf{Classi} & \textbf{Risoluz.} & \textbf{Addestr.} & \textbf{Valid.} & \textbf{Riservati} & \textbf{Durata [s]} \\
\hline
DVS-Lip & 100 & $128\times128$ & $11\,901$ & $2\,995$ & $4\,975$ & $1{,}08$ ($1{,}89$) \\
DVS128 Gesture & 11 & $128\times128$ & 912 & 264 & 288 & $6{,}36$ ($18{,}46$) \\
CIFAR10-DVS & 10 & $128\times128$ & $8\,100$ & 900 & $1\,000$ & $1{,}30$ ($1{,}46$) \\
THU-EACT-50-CHL & 50 & $346\times260$ & $1\,387$ & 399 & 544 & $2{,}71$ ($19{,}91$) \\
\hline
\end{tabular}
\end{table}


% ==============================================================================
%  Rappresentazione dell'input nella baseline
% ==============================================================================

\section{Rappresentazione dell'input}
\label{sec:input_baseline}

La baseline riceve il flusso di eventi come frame di conteggio a finestra fissa (\emph{fixed-window count frames}), la più diretta tra le rappresentazioni introdotte nella sezione~\ref{sec:rappresentazione_event_camera}. A partire dall'origine del campione si considera una finestra di durata $D_{\mathrm{win}}$, la si divide in $T$ bin della stessa durata secondo l'equazione~\eqref{eq:finestre_fisse}, e in ciascun bin si contano, pixel per pixel, gli eventi ON e quelli OFF, come nell'equazione~\eqref{eq:event_frame_accumulo}. Il campione diventa così un tensore
\begin{equation}
\label{eq:input_baseline}
X=\bigl(F_0,\ldots,F_{T-1}\bigr)\in\mathbb N_0^{T\times2\times H\times W},
\qquad
\Delta t=\frac{D_{\mathrm{win}}}{T},
\end{equation}
con un canale per ciascuna polarità. Gli istogrammi di eventi a due canali sono tra le rappresentazioni più diffuse nell'apprendimento profondo su dati neuromorfici~\cite{maqueda2018event,gallego2022event}, e la rete li riceve così come sono, senza alcuna normalizzazione. Nel seguito questa scelta verrà indicata come rappresentazione a finestra fissa.

In tutti i dataset la sequenza conta $T=40$ passi. È un budget comune di passi, non una risoluzione temporale comune, perché la finestra, e con essa la durata dei bin, viene scelta per ciascun dataset in base alle durate di sviluppo (tabella~\ref{tab:finestre_input}). Su tre dataset la finestra contiene per intero anche le registrazioni più lunghe. Su THU-EACT-50-CHL, dove la coda delle durate è molto lunga, una finestra estesa fino al campione più lungo avrebbe richiesto bin di mezzo secondo. Quella di $8\,\mathrm{s}$ conserva invece bin da $200\,\mathrm{ms}$ e il $99\%$ degli eventi, al prezzo di troncare due registrazioni della validazione.

\begin{table}[t]
\centering
\small
\renewcommand{\arraystretch}{1.2}
\caption[Parametri della rappresentazione a finestra fissa]{Parametri della rappresentazione a finestra fissa, con $T=40$ passi in tutti i dataset. Le ultime due colonne si riferiscono alle partizioni di sviluppo. La frazione di finestra occupata è il rapporto tra la durata mediana delle registrazioni e quella della finestra.}
\label{tab:finestre_input}
\begin{tabularx}{\textwidth}{@{}>{\raggedright\arraybackslash}X c c >{\centering\arraybackslash}p{2.0cm} >{\centering\arraybackslash}p{2.0cm}@{}}
\hline
\textbf{Dataset} & \textbf{Finestra} & \textbf{Bin} & \textbf{Eventi conservati} & \textbf{Finestra occupata} \\
\hline
DVS-Lip & $2\,\mathrm{s}$ & $50\,\mathrm{ms}$ & $100\%$ & $54\%$ \\
DVS128 Gesture & $20\,\mathrm{s}$ & $500\,\mathrm{ms}$ & $100\%$ & $32\%$ \\
CIFAR10-DVS & $2\,\mathrm{s}$ & $50\,\mathrm{ms}$ & $100\%$ & $65\%$ \\
THU-EACT-50-CHL & $8\,\mathrm{s}$ & $200\,\mathrm{ms}$ & $99{,}06\%$ & $34\%$ \\
\hline
\end{tabularx}
\end{table}

Una finestra prefissata ha però un effetto collaterale, ben visibile nella figura~\ref{fig:esempi_dataset}. Le registrazioni più brevi della finestra ne lasciano vuoti gli ultimi bin, che non sono padding ma osservazioni valide di un intervallo senza attività (sezione~\ref{subsec:tempo_valido_padding}). La rete li elabora comunque, e il readout ne include i contributi nella propria media. Su DVS128 Gesture e THU-EACT-50-CHL un campione tipico trascorre così circa due terzi dei passi senza ricevere alcun evento. Il capitolo~\ref{chap:clock_input} analizza le conseguenze di questa scelta e la confronta con altri clock. Fino ad allora la rappresentazione a finestra fissa resta invariata, perché ogni differenza osservata nel capitolo~\ref{chap:riallocazione_budget} possa essere attribuita all'architettura.


% ==============================================================================
%  Architettura della baseline compatta
% ==============================================================================

\section{Architettura della baseline compatta}
\label{sec:architettura_baseline}

La baseline compatta è un'implementazione scritta per questo lavoro, che conserva i principi di QKFormer (sezione~\ref{sec:qkformer_riferimento}) riducendone drasticamente la scala. Il backbone ha due stadi gerarchici (figura~\ref{fig:architettura_baseline}a). Un primo patch embedding porta l'ingresso a un ottavo della risoluzione con 64 canali, un blocco con Q-K Token Attention elabora questa rappresentazione, un secondo patch embedding dimezza ancora la risoluzione portando i canali a 128, e un blocco con spiking self-attention precede il readout. Le differenze rispetto all'architettura originale sono riassunte alla fine della sezione.

\begin{figure}[t]
    \centering
    \includegraphics[width=\textwidth]{images/cap2/cap2_architettura.pdf}
    \caption[Architettura della baseline compatta]{Architettura della baseline compatta per un ingresso di $128\times128$ pixel. (a)~Sequenza dei blocchi, con le dimensioni di uscita di ciascuno per ogni passo, nella forma canali $\times$ altezza $\times$ larghezza. (b)~Il primo patch embedding, l'unico modulo che lavora a piena risoluzione, con il ramo principale di quattro convoluzioni e la scorciatoia che parte dalla prima mappa di feature. (c)~Ripartizione dei parametri e dello stato persistente tra i blocchi su DVS-Lip, secondo la tabella~\ref{tab:architettura_baseline}. I colori dei blocchi corrispondono a quelli del pannello (a).}
    \label{fig:architettura_baseline}
\end{figure}

\paragraph{Neurone e convenzioni di calcolo}
\label{subsec:neurone_convenzioni}

Tutti i neuroni della rete sono LIF, discretizzati con lo schema di Eulero dell'equazione~\eqref{eq:lif_euler_discreto} con $\tau_d=2$ e dotati del soft reset dell'equazione~\eqref{eq:soft_reset} con soglia $V_{\mathrm{th}}=1$. Il potenziale viene azzerato all'inizio di ogni campione. Fanno eccezione i neuroni che prendono le decisioni dell'attenzione, descritti più avanti, la cui soglia è dimezzata. In addestramento la derivata dello spike è sostituita dalla surrogata sigmoidale dell'equazione~\eqref{eq:surrogate_sigmoide} con $\alpha=4$, e il termine di reset è escluso dal calcolo del gradiente.

Con $\tau_d=2$ il potenziale si dimezza a ogni passo, e la costante di tempo equivalente dell'equazione~\eqref{eq:tau_fisica_equivalente} vale circa $1{,}44\,\Delta t$. Lo stesso neurone dimentica quindi in tempi fisici molto diversi da un dataset all'altro, in circa $72\,\mathrm{ms}$ su DVS-Lip e CIFAR10-DVS, $290\,\mathrm{ms}$ su THU-EACT-50-CHL e $720\,\mathrm{ms}$ su DVS128 Gesture. È un primo segno di quanto il clock dell'input fissi anche l'orizzonte della memoria neuronale, un tema che il capitolo~\ref{chap:clock_input} riprende.

Convoluzioni, normalizzazioni e proiezioni lineari vengono applicate a ogni passo in modo indipendente, come se i passi fossero campioni distinti. L'unico canale attraverso cui l'informazione passa da un istante al successivo è quindi lo stato dei neuroni, come in QKFormer (sezione~\ref{subsec:qkformer_tempo}). Anche i collegamenti residui seguono QKFormer (sezione~\ref{subsec:qk_attention}). Sommano spike a valle dei neuroni, e i tensori scambiati tra i blocchi contengono perciò piccoli interi anziché valori binari, con conseguenze dirette sul conteggio delle operazioni (sezione~\ref{sec:caratterizzazione_computazionale}).

\paragraph{Primo patch embedding}
\label{subsec:primo_embedding}
Il primo patch embedding traduce i conteggi in spike e riduce la risoluzione di un fattore otto (figura~\ref{fig:architettura_baseline}b). Il ramo principale è una cascata di quattro convoluzioni $3\times3$, ciascuna seguita da batch normalization e da un neurone LIF. La prima porta i canali da 2 a 32 a piena risoluzione, le altre tre dimezzano a turno la risoluzione e portano i canali a 64. In parallelo, una scorciatoia preleva le feature del primo strato e le porta direttamente alla risoluzione finale con una convoluzione $1\times1$ di passo 8. Le due uscite, entrambe spiking, vengono sommate,
\begin{equation}
\label{eq:primo_embedding}
E^{(1)}_k=M^{(1)}_k+R^{(1)}_k\in\{0,1,2\}^{64\times\frac H8\times\frac W8},
\end{equation}
dove $M^{(1)}_k$ e $R^{(1)}_k$ sono gli spike emessi al passo $k$ dal ramo principale e dalla scorciatoia. È il principio della scorciatoia deformata di SPEDS (equazione~\eqref{eq:speds_formulazione}), con due differenze. Il sottocampionamento avviene con convoluzioni a passo 2 anziché con max-pooling, e la scorciatoia attraversa l'intero cambio di scala in un solo salto. È anche l'unico modulo della rete che opera a piena risoluzione, e il suo peso sul bilancio complessivo emerge più avanti.

\paragraph{Stadio con Q-K Token Attention}
\label{subsec:stadio_qkta}
Il primo stadio contiene un unico blocco Transformer con Q-K Token Attention. Per un ingresso di $128\times128$ pixel lavora su $16\times16=256$ token con 64 canali, divisi in 8 teste da 8 canali. Query e Key si ottengono come nell'equazione~\eqref{eq:qk_projections_qkformer}, e per ogni testa $h$ e ogni token $n$ la somma degli spike della Query sui canali della testa attraversa un neurone con soglia $0{,}5$,
\begin{equation}
\label{eq:qkta_baseline}
a_{h,n}=\operatorname{SN}_{0{,}5}\Bigl(\sum_{d\in h}Q_{n,d}\Bigr),
\qquad
X'_{n,d}=a_{h,n}\,K_{n,d}\quad\text{per } d\in h.
\end{equation}
La decisione binaria lascia passare o azzera i canali della Key di quella testa per quel token. Rispetto all'equazione~\eqref{eq:qkta_vettore_attenzione} ogni token dispone così di otto decisioni indipendenti, una per testa. Il neurone che le produce possiede inoltre uno stato, e la selezione di un token dipende anche dalla storia recente della sua Query. Seguono la proiezione di uscita dell'equazione~\eqref{eq:qk_proiezione_lineare_uscita}, il collegamento residuo e uno spiking MLP che espande i canali da 64 a 128 e li riporta a 64, secondo lo schema delle equazioni~\eqref{eq:blocco_qkformer_qka} e~\eqref{eq:blocco_qkformer_smlp}.

\paragraph{Secondo patch embedding e stadio con SSA}
\label{subsec:stadio_ssa}
Il secondo patch embedding segue lo stesso schema a due rami. Il ramo principale applica una convoluzione $3\times3$ da 64 a 128 canali e una seconda con passo 2, la scorciatoia una convoluzione $1\times1$ con passo 2. Ne risultano 128 canali a un sedicesimo della risoluzione d'ingresso, cioè 64 token. Il secondo stadio impiega la SSA, con Query, Key e Value spiking~\eqref{eq:spiking_qkv} divise in 8 teste da 16 canali. Per ciascuna testa
\begin{equation}
\label{eq:ssa_baseline}
Y=\operatorname{SN}_{0{,}5}\bigl(\beta\,Q(K^{\top}V)\bigr),
\qquad
\beta=\frac{1}{\sqrt{16}},
\end{equation}
cui seguono la proiezione di uscita, il collegamento residuo e uno spiking MLP che porta i canali da 128 a 256 e di nuovo a 128. In assenza di softmax il prodotto è associativo, e l'implementazione calcola prima $K^{\top}V$, una matrice $16\times16$ per testa. Il risultato coincide con quello dell'equazione~\eqref{eq:ssa_spikformer}, ma la matrice $N\times N$ tra token non viene mai costruita e il costo cresce linearmente con il numero dei token.

\paragraph{Readout}
\label{subsec:readout_baseline}
Il readout media la rappresentazione finale prima sulle posizioni e poi sui passi,
\begin{equation}
\label{eq:readout_media}
z=\frac{1}{T\,N_2}\sum_{k=0}^{T-1}\sum_{n=1}^{N_2}Z^{(2)}_{k,n},
\qquad
\hat y=W_c\,z+b_c,
\end{equation}
dove $Z^{(2)}_{k,n}\in\mathbb R^{128}$ è l'uscita del secondo stadio al passo $k$ e nella posizione $n$, $N_2$ il numero dei suoi token, $W_c$ e $b_c$ i parametri della testa lineare. Le feature mediate appartengono al flusso residuo e non sono binarie. La media comprende tutti i $T$ passi, anche quelli successivi all'ultimo evento, e la decisione riassume così l'intera sequenza, con lo stesso peso per ogni passo.

\paragraph{Bilancio di parametri e stato}
\label{subsec:bilancio_baseline}

La tabella~\ref{tab:architettura_baseline} riporta, blocco per blocco, le dimensioni, i parametri e lo stato persistente, cioè il numero di potenziali di membrana che devono sopravvivere da un passo al successivo. Il backbone conta $487\,808$ parametri, ai quali si aggiunge la testa, che dipende dal numero delle classi, per un totale che va da $489\,098$ su CIFAR10-DVS a $500\,708$ su DVS-Lip. Su THU-EACT-50-CHL, con un ingresso di $346\times260$ pixel, i due stadi lavorano su $1\,452$ e 374 token, e lo stato cresce in proporzione al numero delle posizioni.

\begin{table}[t]
\centering
\small
\renewcommand{\arraystretch}{1.2}
\caption[Struttura della baseline compatta]{Struttura della baseline compatta per un ingresso di $128\times128$ pixel. Le uscite sono indicate come canali $\times$ altezza $\times$ larghezza per ciascun passo. Lo stato conta i potenziali di membrana di tutti i neuroni del blocco. I parametri della testa dipendono dal numero $C$ delle classi, e il totale si riferisce a DVS-Lip.}
\label{tab:architettura_baseline}
\begin{tabularx}{\textwidth}{@{}>{\raggedright\arraybackslash}p{2.6cm} >{\raggedright\arraybackslash}X c r r@{}}
\hline
\textbf{Blocco} & \textbf{Operazioni principali} & \textbf{Uscita} & \textbf{Parametri} & \textbf{Stato} \\
\hline
Primo patch embedding & 4 conv.\ $3\times3$ ($2\to32\to64\to64\to64$), scorciatoia $1\times1$ a passo 8 & $64\times16\times16$ & $95\,360$ & $884\,736$ \\
Stadio 1 & QKTA a 8 teste, MLP $64\to128\to64$ & $64\times16\times16$ & $29\,440$ & $100\,352$ \\
Secondo patch embedding & conv.\ $3\times3$ ($64\to128$), conv.\ $3\times3$ a passo 2, scorciatoia $1\times1$ a passo 2 & $128\times8\times8$ & $230\,144$ & $49\,152$ \\
Stadio 2 & SSA a 8 teste, MLP $128\to256\to128$ & $128\times8\times8$ & $132\,864$ & $65\,536$ \\
Readout e testa & media su spazio e tempo, lineare $128\to C$ & $C$ & $129\,C$ & 0 \\
\hline
\textbf{Totale} & DVS-Lip, $C=100$ & & $500\,708$ & $1\,099\,776$ \\
\hline
\end{tabularx}
\end{table}

La figura~\ref{fig:architettura_baseline}c mostra che parametri e stato si distribuiscono in modo quasi speculare. I due patch embedding contengono circa due terzi dei parametri, e il secondo da solo quasi la metà. Lo stato è invece dominato dal primo, che lavora sulle mappe più risolute e ne ospita oltre l'$80\%$, e nel complesso supera di più del doppio il numero dei parametri. Nel regime compatto capacità e costo si concentrano dunque lontano dall'attenzione, ed è da questa osservazione che parte il capitolo~\ref{chap:riallocazione_budget}.

\paragraph{Differenze rispetto a QKFormer}
\label{subsec:differenze_qkformer}

Rispetto alla configurazione neuromorfica di QKFormer (sezione~\ref{subsec:gerarchia_patch}), la baseline conserva la sequenza di un blocco QKTA e di un blocco SSA, il patch embedding con scorciatoia e la separazione tra memoria neuronale e interazioni all'interno del passo. Cambiano la scala e alcuni dettagli costruttivi, riassunti nella tabella~\ref{tab:confronto_qkformer}. La differenza nel numero degli stadi è meno marcata di quanto sembri, perché nella configurazione originale il primo stadio non contiene blocchi e si riduce al patch embedding. Più sostanziali sono la larghezza massima, dimezzata da 256 a 128 canali, il numero dei parametri, sceso a circa un terzo, e la lunghezza delle sequenze, di 40 passi anziché 16. L'implementazione, infine, mantiene il tempo come asse separato prima del neurone di uscita della SSA, il punto in cui gli autori hanno segnalato un errore del codice ufficiale~\cite{qkformerImplementationNotice}, e nella baseline quel problema non si presenta.

\begin{table}[t]
\centering
\small
\renewcommand{\arraystretch}{1.2}
\caption[QKFormer neuromorfico e baseline compatta a confronto]{Confronto tra la configurazione neuromorfica di QKFormer, secondo la descrizione della sezione~\ref{sec:qkformer_riferimento}, e la baseline compatta.}
\label{tab:confronto_qkformer}
\begin{tabularx}{\textwidth}{@{}>{\raggedright\arraybackslash}p{3.8cm} >{\raggedright\arraybackslash}X >{\raggedright\arraybackslash}X@{}}
\hline
\textbf{Aspetto} & \textbf{QKFormer, configurazione neuromorfica} & \textbf{Baseline compatta} \\
\hline
Stadi e blocchi & tre stadi, blocchi $\{0,1,1\}$ & due stadi, un blocco ciascuno \\
Attenzione & QKTA, poi SSA & QKTA, poi SSA \\
Larghezza massima & 256 canali & 128 canali \\
Parametri & circa $1{,}50$ milioni & da $489\,098$ a $500\,708$ \\
Sottocampionamento & convoluzione, BN, neurone e max-pooling & convoluzioni a passo 2 \\
Passi temporali & 16 & 40 \\
\hline
\end{tabularx}
\end{table}


% ==============================================================================
%  Protocollo sperimentale
% ==============================================================================

\section{Protocollo sperimentale}
\label{sec:protocollo_addestramento}

Tutti i modelli vengono addestrati da zero con la stessa ricetta, e due modelli messi a confronto differiscono per un solo elemento. Questa sezione descrive la ricetta, i vincoli che ne derivano per i confronti e i modi in cui un componente della rete può essere interrogato.

\paragraph{Ricetta di addestramento}
\label{subsec:ricetta}
La ricetta, riassunta nella tabella~\ref{tab:ricetta_addestramento}, è comune a tutti i modelli e a tutti i dataset. L'ottimizzatore è AdamW~\cite{loshchilov2019adamw}, con un tasso di apprendimento che cresce nelle prime epoche e decresce poi con un profilo cosinusoidale~\cite{loshchilov2017sgdr} lungo 128 epoche. La funzione di perdita è l'entropia incrociata con label smoothing~\cite{szegedy2016rethinking}, la norma del gradiente viene limitata~\cite{pascanu2013difficulty} e i calcoli avvengono in precisione mista~\cite{micikevicius2018mixed}. L'unica augmentation è un ribaltamento orizzontale su DVS-Lip. Sugli altri dataset non si applica alcuna trasformazione, e su DVS128 Gesture il ribaltamento sarebbe anche scorretto, perché scambierebbe tra loro classi direzionali. La ricetta non viene adattata ai singoli modelli né ai singoli dataset, e quando una campagna la modifica, la modifica vale per tutti i modelli del confronto.

\paragraph{Microbatch e batch normalization}
\label{subsec:microbatch_bn}
Il batch effettivo di 32 campioni si ottiene accumulando i gradienti di più microbatch, la cui dimensione dipende dalla memoria richiesta da ciascun dataset. L'accumulo rende equivalente il gradiente delle componenti lineari, ma non il comportamento della batch normalization~\cite{ioffe2015batch}, che in addestramento normalizza con le statistiche del solo microbatch corrente. Poiché il tempo è fuso con la dimensione del batch, quelle statistiche aggregano tutti i passi e tutti i campioni del microbatch, come avviene comunemente nelle SNN addestrate direttamente~\cite{zheng2021going}. Ne discendono due conseguenze. Il frazionamento in microbatch fa parte del protocollo, perché quattro microbatch da 8 campioni non equivalgono a un batch da 32, e ogni eccezione viene dichiarata dove si presenta. Inoltre, in addestramento ogni passo viene normalizzato con statistiche che dipendono anche dai passi successivi, e la rete non è strettamente causale, pur essendolo ciascuno dei suoi operatori. In inferenza le statistiche sono fisse e ogni passo torna a dipendere soltanto da quelli precedenti.

\begin{table}[t]
\centering
\small
\renewcommand{\arraystretch}{1.2}
\caption[Ricetta di addestramento comune]{Ricetta di addestramento comune a tutti i modelli e a tutti i dataset.}
\label{tab:ricetta_addestramento}
\begin{tabularx}{\textwidth}{@{}>{\raggedright\arraybackslash}p{3.6cm} >{\raggedright\arraybackslash}X@{}}
\hline
\textbf{Componente} & \textbf{Scelta} \\
\hline
Ottimizzatore & AdamW, tasso di apprendimento $3\times10^{-4}$, weight decay $5\times10^{-4}$ \\
Schedule & warmup lineare di 4 epoche da $0{,}01$ volte il valore nominale, poi decadimento cosinusoidale fino a $10^{-6}$ \\
Durata & 128 epoche \\
Funzione di perdita & entropia incrociata con label smoothing $0{,}1$ \\
Stabilizzazione & limite di $1$ sulla norma globale del gradiente \\
Precisione & precisione mista automatica su GPU \\
Batch effettivo & 32 campioni, ottenuti per accumulo su più microbatch \\
Augmentation & ribaltamento orizzontale con probabilità $0{,}5$, solo su DVS-Lip \\
Selezione del checkpoint & epoca con la massima Macro-F1 sulla validazione di sviluppo \\
\hline
\end{tabularx}
\end{table}

\paragraph{Test di sovradattamento}
\label{subsec:verifiche_preliminari}
Prima di un addestramento completo, ogni nuova configurazione deve superare un test di sovradattamento controllato. Il modello viene addestrato su pochi campioni, quattro per ciascuna di al più 16 classi, e deve imparare a classificarli quasi perfettamente, mantenendo per cinque epoche consecutive un'accuratezza superiore al $95\%$. Il test non dice nulla sulla capacità di generalizzare, ma mostra che il modello è in grado di apprendere e intercetta a un costo trascurabile errori che altrimenti emergerebbero soltanto alla fine di un addestramento completo. Una configurazione che non lo supera non viene addestrata.

\paragraph{Confronti tra modelli}
\label{sec:disegno_controlli}
Due modelli messi a confronto differiscono per un solo elemento, mentre dati, partizioni, ingresso, ricetta e regola di selezione restano identici. Quando l'oggetto dell'esperimento è l'ingresso, è l'architettura a restare ferma. Lo stesso componente, però, può essere interrogato in tre modi, che rispondono a domande diverse.

Un'ablazione riaddestrata rimuove o sostituisce il componente e addestra il modello da zero. Misura che cosa riesca a ottenere l'architettura senza di esso, ed è la sola prova che il componente serva davvero.

Un intervento sul checkpoint modifica invece un modello già addestrato, soltanto in inferenza. Spegnere un ramo, eliminare un ritardo o rimescolare l'ordine dei passi misura quanto la soluzione appresa dipenda da quel componente, non se il componente sia necessario. Una rete addestrata con un modulo può essersi adattata alla sua presenza e degradare quando lo si rimuove, anche se un'architettura che ne è priva avrebbe raggiunto da zero un risultato analogo. Il ruolo dell'attenzione nel secondo stadio viene esaminato combinando i due punti di vista, con un test di necessità che la disattiva sul checkpoint addestrato e un test di sostituibilità che la rimpiazza con un mixer spaziale locale e riaddestra il modello.

Un controllo di capacità, infine, separa l'effetto di un operatore da quello delle risorse che aggiunge. Ogni nuovo modulo porta con sé parametri e operazioni, e un miglioramento potrebbe dipendere soltanto da questi. Il confronto con un'alternativa di costo comparabile ma priva della struttura in esame, come un allargamento uniforme della rete o un percettrone multistrato con lo stesso numero di parametri, permette di attribuire il guadagno alla struttura e non alla quantità.


% ==============================================================================
%  Valutazione e incertezza
% ==============================================================================

\section{Valutazione e incertezza}
\label{sec:metriche_valutazione}

\paragraph{Metriche}
\label{subsec:macro_f1}
La metrica principale è la Macro-F1, cioè la media non pesata, sulle $C$ classi, degli F1 delle singole classi~\cite{sokolova2009systematic},
\begin{equation}
\label{eq:macro_f1}
\text{Macro-F1}=\frac{1}{C}\sum_{c=1}^{C}\frac{2P_cR_c}{P_c+R_c},
\end{equation}
dove $P_c$ e $R_c$ sono la precisione e il richiamo della classe $c$. A differenza dell'accuratezza, la Macro-F1 attribuisce a tutte le classi lo stesso peso, indipendentemente dalla loro numerosità, e penalizza un modello che ottenga un buon risultato complessivo trascurando alcune classi.\footnote{Su THU-EACT-50-CHL la media comprende tutte le 50 classi della testa, anche se quattro di esse compaiono soltanto in soggetti del training e mancano sia nella validazione di sviluppo sia nel test ufficiale. Il loro F1 è nullo per costruzione, e il valore massimo raggiungibile della Macro-F1 è del $92\%$. Il vincolo è identico per tutti i modelli e non altera i confronti, ma va tenuto presente nella lettura dei valori assoluti.} Accanto a essa si riportano l'accuratezza e, su DVS-Lip, le accuratezze Acc1 e Acc2 dei due gruppi di parole descritti nella sezione~\ref{sec:dataset}.

\paragraph{Selezione del checkpoint}
\label{subsec:selezione_incertezza}
Ogni modello viene valutato sulla validazione di sviluppo al termine di ciascuna epoca, e il checkpoint conservato è quello con la Macro-F1 più alta. Selezione e misura avvengono quindi sugli stessi dati, e il valore riportato è una stima ottimistica della prestazione su dati nuovi, tanto più quanto più la curva di validazione oscilla. Per questo si registrano anche la media e la deviazione standard della Macro-F1 nelle ultime 16 epoche, un indicatore di stabilità meno esposto alla scelta del massimo.

\paragraph{Incertezza dei confronti}
Quando due modelli vengono valutati sulla stessa validazione, le loro predizioni sono appaiate campione per campione. Il test esatto di McNemar considera soltanto i campioni su cui i due modelli discordano, e verifica se gli errori corretti dall'uno e quelli introdotti dall'altro siano sbilanciati~\cite{mcnemar1947note,dietterich1998approximate}. Il bootstrap appaiato ricampiona la validazione e ricava un intervallo di confidenza al $95\%$ per la differenza di accuratezza~\cite{efron1993bootstrap}. I valori di $p$ non sono corretti per i confronti multipli, e vanno letti come indicazioni più che come verdetti.

Entrambi gli strumenti misurano l'incertezza dovuta alla dimensione finita della validazione, a modelli fissati. Nessuno dei due misura la variabilità dell'addestramento, cioè quanto cambierebbe il risultato ripetendolo con un'inizializzazione diversa, che può essere dello stesso ordine delle differenze tra i modelli confrontati~\cite{bouthillier2021accounting}. Le esplorazioni usano perciò un solo addestramento per configurazione, e i loro risultati sono trattati come indicazioni. Le configurazioni destinate a diventare un riferimento vengono ripetute tre volte, e solo per queste si riportano media e deviazione standard.

\paragraph{Valutazione sui prefissi}
\label{subsec:prefissi}
Su DVS-Lip e DVS128 Gesture si misura anche come la decisione si formi nel tempo. Lo stesso checkpoint viene valutato su porzioni iniziali della sequenza, con il readout calcolato sui soli passi osservati, e un indice sintetico, indicato come PrefixAUC, riassume l'andamento dell'accuratezza al crescere della porzione osservata. È una misura diagnostica. Non interviene nella selezione dei modelli e non autorizza a ridurre la finestra di ingresso, perché un modello addestrato su sequenze complete non è stato ottimizzato per decidere su sequenze troncate.

\paragraph{Partizioni riservate}
\label{subsec:partizioni_riservate}
Tutte le metriche riportate nei capitoli successivi sono misurate sulla validazione di sviluppo, salvo indicazione esplicita. Le partizioni riservate restano escluse dallo sviluppo, dalla selezione dei checkpoint e dalla caratterizzazione computazionale. L'unica eccezione riguarda alcune analisi diagnostiche, presentate separatamente, in cui checkpoint già selezionati su DVS-Lip e THU-EACT-50-CHL sono stati valutati in sola inferenza sul test ufficiale. Quei risultati descrivono il comportamento di modelli già scelti e non vengono usati per sceglierne altri. % TODO: \ref al capitolo dell'analisi trasversale (diagnostica sul test ufficiale).


% ==============================================================================
%  Caratterizzazione computazionale
% ==============================================================================

\section{Caratterizzazione computazionale}
\label{sec:caratterizzazione_computazionale}

Il costo di ciascun modello viene misurato in inferenza sul checkpoint selezionato, usando un sottoinsieme fisso di 64 campioni della validazione di sviluppo, lo stesso per tutti i modelli di un dataset. Le grandezze raccolte sono quelle introdotte nella sezione~\ref{sec:efficienza_neuromorfica_vincoli}, e la tabella~\ref{tab:contatori_profilo} ne precisa il significato.

\begin{table}[t]
\centering
\small
\renewcommand{\arraystretch}{1.2}
\caption[Grandezze della caratterizzazione computazionale]{Grandezze raccolte nella caratterizzazione computazionale di ciascun checkpoint.}
\label{tab:contatori_profilo}
\begin{tabularx}{\textwidth}{@{}>{\raggedright\arraybackslash}p{3.2cm} >{\raggedright\arraybackslash}X@{}}
\hline
\textbf{Grandezza} & \textbf{Che cosa conta} \\
\hline
Parametri & pesi addestrabili, testa inclusa \\
Stato persistente & potenziali di membrana e buffer temporali che devono sopravvivere da un passo al successivo \\
MAC & moltiplicazioni e accumuli con un operando non binario, come i conteggi in ingresso, le somme residue e le correnti dei moduli temporali \\
AC potenziali & accumuli con ingresso binario, contati come se ogni possibile spike fosse presente \\
SOP & accumuli effettivamente innescati dagli spike osservati, comprese le operazioni dell'attenzione, secondo l'equazione~\eqref{eq:stima_sop} \\
Firing rate & frazione di spike emessi da ciascuno strato \\
\hline
\end{tabularx}
\end{table}

La distinzione tra MAC e accumuli segue il tipo di operando, non il tipo di strato. Uno strato che legge spike esegue accumuli, mentre lo stesso strato applicato a conteggi o a somme residue esegue moltiplicazioni vere e proprie. Per questo i collegamenti residui della baseline, che producono piccoli interi (sezione~\ref{sec:architettura_baseline}), pesano sul bilancio, perché una parte delle proiezioni della rete riceve valori non binari.

Per riassumere i conteggi in un'unica grandezza si adotta la stima dell'equazione~\eqref{eq:energia_snn_semplificata}, con i costi di riferimento per operazioni a 32 bit in tecnologia a $45\,\mathrm{nm}$ riportati da Horowitz~\cite{horowitz2014energy}, cioè $4{,}6\,\mathrm{pJ}$ per una MAC, $0{,}9\,\mathrm{pJ}$ per un'addizione e $3{,}7\,\mathrm{pJ}$ per una moltiplicazione. La stima viene calcolata in due varianti, una densa, basata sugli accumuli potenziali, e una pesata per l'attività osservata. Come discusso nella sezione~\ref{sec:efficienza_neuromorfica_vincoli}, è una proxy parziale del solo costo aritmetico, utile per confrontare architetture a parità di input e in nessun caso una misura dell'energia consumata. Lo stato persistente, infine, viene riportato sempre accanto ai parametri, perché nella baseline le due grandezze differiscono di più di un fattore due e un modulo con pochi pesi può richiedere un buffer molto più grande.


% ==============================================================================
%  Il riferimento nei capitoli successivi
% ==============================================================================

\section{Il riferimento nei capitoli successivi}
\label{sec:sintesi_protocollo}

Dataset e partizioni, ricetta di addestramento, regola di selezione del checkpoint, metriche e grandezze di costo restano gli stessi in tutti i capitoli sperimentali. Cambiano invece, da un capitolo all'altro, la rete di partenza e l'ingresso, secondo il percorso descritto all'inizio del capitolo e riassunto nella tabella~\ref{tab:mappa_riferimenti}. In ogni capitolo varia un solo fattore, e i confronti avvengono sempre con il riferimento di quel capitolo. I nomi con cui modelli e varianti compaiono nel testo sono descrittivi, e la loro corrispondenza con le esecuzioni effettive è riportata in appendice. % TODO: \ref all'appendice di riproducibilità

\par\medskip\noindent
\begin{minipage}{\textwidth}
\small
\renewcommand{\arraystretch}{1.2}
\captionof{table}[Il riferimento nei capitoli sperimentali]{Rete di partenza, ingresso e fattore studiato nei capitoli sperimentali. Tutti gli altri elementi del protocollo restano quelli descritti in questo capitolo.}
\label{tab:mappa_riferimenti}
\begin{tabularx}{\textwidth}{@{}>{\raggedright\arraybackslash}p{3.3cm} >{\raggedright\arraybackslash}p{3.0cm} >{\raggedright\arraybackslash}X >{\raggedright\arraybackslash}X@{}}
\hline
\textbf{Capitolo} & \textbf{Rete di partenza} & \textbf{Ingresso} & \textbf{Fattore studiato} \\
\hline
\ref{chap:riallocazione_budget}, riallocazione del budget & baseline compatta & finestra fissa & distribuzione della capacità nella rete \\
\ref{chap:clock_input}, clock dell'input & modello di riferimento & finestra fissa, eventi costanti, lunghezza variabile & rappresentazione dell'ingresso \\
\ref{chap:memoria_relazionale}, memoria relazionale & modello di riferimento & lunghezza variabile & memoria temporale \\
\hline
\end{tabularx}
\end{minipage}
\par\medskip

Il capitolo successivo parte dal bilancio della sezione~\ref{sec:architettura_baseline}, che concentra stato e operazioni lontano dall'attenzione, e si chiede come spendere meglio un budget di capacità così ridotto.
